\section{W[1]-hardness}\label{sec:hardness}
Theorems~\ref{thm:l0} and~\ref{thm:A} solve CEaMP in time $5^{w+1}n^{2\ell+O(1)}$, and Theorem~\ref{thm:nl} halves the exponent of $n$; all of these are XP in $\ell$. This section shows that the exponent of $n$ cannot be removed: the main question of this paper has a negative answer.
\begin{theorem}[{W[1]}-hardness]\label{thm:hard}
There is a polynomial-time reduction that maps an instance $(k,n,(S_{ij}))$ of Grid Tiling with $\le$ to an instance $(G,H,\ell)$ of CEaMP with the following properties. Let $p$ be the least even integer with $p\ge k^2+5$.
\begin{enumerate}
\item[(i)] $\ell=k^2$; $H$ is a pair-maximal modification-disjoint packing; $|V(G)|=O(k^2n^2p)$.
\item[(ii)] $c(H)\le 4p+2$, and $\tw(I(H))\le\mathrm{pw}(I(H))\le 24\,k\,p^2$, independently of $n$; $I(H)$ has at most $2k$ components.
\item[(iii)] $(G,H,\ell)$ is a YES instance iff the Grid Tiling instance has a solution; if it has one, then $\OPT(G)=|H|+k^2$.
\end{enumerate}
Consequently CEaMP is W[1]-hard parameterized by $\ell+\tw(I(H))$, by $\ell+\mathrm{pw}(I(H))$ and by $\ell+c$, even for pair-maximal packings. Unless $\mathrm{FPT}=\mathrm{W[1]}$ there is no algorithm with running time $f(\ell,w)\,n^{O(1)}$, and unless ETH fails there is no algorithm with running time $f(\ell,w)\,N^{o(\sqrt{\ell})}$, where $N=|V(G)|$.
\end{theorem}
Here $p\le k^2+6$, so on the produced instances $\ell=k^2$, $c\le 4k^2+26$ and $\tw(I(H))\le 24k(k^2+6)^2$ are all functions of $k$ alone. Grid Tiling with $\le$ is the following problem~\cite{Cygan15} (Chapter~14): given $k$, $n$ and non-empty sets $S_{ij}\subseteq[n]^2$ ($1\le i,j\le k$), choose $s_{ij}=(a_{ij},b_{ij})\in S_{ij}$ with $a_{ij}\le a_{i,j+1}$ and $b_{ij}\le b_{i+1,j}$; it is W[1]-hard in $k$ and has no $f(k)n^{o(k)}$ algorithm under ETH (Marx; see~\cite{Cygan15}, Chapter~14). (We use the version of the cited chapter after exchanging the two coordinates.) The reduction, all lemmas and the proof of Theorem~\ref{thm:hard} are in Appendix~\ref{app:hardness}; this section states what is proved, outlines the construction, and says precisely which steps of the proof use a computer and what has been tested.

\subsection{What the theorem says, and what it does not}
\begin{corollary}[The ETH bound that follows]\label{cor:hardeth}
Assume ETH. There are no computable $f$ and no algorithm that decides CEaMP in time $f(\ell,\tw(I(H)))\,N^{o(\sqrt\ell)}$ on instances with $N$ vertices.
\end{corollary}
\begin{proof}
Suppose such an algorithm runs in time $f(\ell,w)N^{h(\ell)}$ with $h(\ell)=o(\sqrt\ell)$. Apply the reduction of Theorem~\ref{thm:hard} to a Grid Tiling instance with parameter $k$: then $\ell=k^2$, $w\le 24k(k^2+6)^2=:g(k)$ and $N\le \mathrm{poly}(k)\,n^2$. The running time is at most $f(k^2,g(k))\,(\mathrm{poly}(k)n^2)^{h(k^2)}$, and $h(k^2)=o(k)$, so it is at most $f'(k)\,n^{o(k)}$ for a computable $f'$ (the factor $\mathrm{poly}(k)^{h(k^2)}=2^{o(k\log k)}$ is absorbed in $f'$). This solves Grid Tiling with $\le$ in time $f'(k)n^{o(k)}$, contradicting ETH~\cite{Cygan15}.
\end{proof}
Nothing stronger follows from this reduction: since $\ell=k^2$, the lower bound is $n^{o(\sqrt\ell)}$, while the upper bounds of Theorems~\ref{thm:A} and~\ref{thm:nl} have exponent $2\ell$ and $\ell$. The gap between $\sqrt\ell$ and $\ell$ is open (Open Problem~\ref{op:eth}). Since $w$ is a polynomial in $k$ on the produced instances, no statement in terms of $\ell+w$ jointly is claimed beyond Corollary~\ref{cor:hardeth}.

\emph{Matching the XP algorithm.} Theorem~\ref{thm:hard} shows that the exponent of $n$ in Theorems~\ref{thm:A} and~\ref{thm:nl} cannot be made independent of $\ell$, even though $w$ and $c$ are bounded by functions of $\ell$ alone on the hard instances. Up to the exact exponent of $n$ ($\Omega(\sqrt\ell)$ under ETH, against $\ell$ resp.\ $2\ell$ proved), the complexity of the pair $(\ell,w)$ is therefore settled.

\emph{What is still open.} (1) The reduction has $c=4p+2$ with $p\ge k^2+5$, so $c$ grows with $\ell$; whether CEaMP is W[1]-hard (or fixed-parameter tractable in $\ell$) when $c$ is bounded by a constant, in particular for $c=2$, is open (Open Problems~\ref{op:boundc} and~\ref{op:c2}). The reason is local to the proof: $p$ is the size of the cliques into which every chain vertex is blown up, and the argument that no clique is split (Lemma~\ref{lem:hsplit}) needs $p>\ell$. (2) The exact exponent under ETH (Open Problem~\ref{op:eth}). (3) Packings with $w\le1$ with tree components of unbounded size, and soft-triangle-free packings (Open Problems~\ref{op:forest} and~\ref{op:BC}); the instances here have many soft triangles and large $w$ (the width is $\Theta(p^2)$ per chain).

\subsection{The construction in outline}
The reduction follows the plan of the tunnel attempt of Appendix~\ref{app:tunnels}, and realises the abstract problem of Theorem~\ref{thm:B} with soft triangles; the obstruction of Theorem~\ref{thm:C} is circumvented because the instances are not soft-triangle-free. We use the language of zero-excess (ZE) solutions: a member $(a,b,c)$ with centre $b$ is in state $L$, $R$ or $M$ if its single edit deletes $ab$, deletes $bc$, or inserts $ac$ (states $1,2,3$ of Appendix~\ref{app:tunnels}).

\paragraph{A selection chain.} A \emph{thin chain} $\mathrm{OH}_m$ ($m\ge7$) is an explicit tight instance on the vertices $s_0,\dots,s_{2m},e_0,\dots,e_m$ with spine members $P_j=(s_{2j},s_{2j+1},s_{2j+2})$, further interior members $G_j,F_j$ that carry the vertices $e_j$, two caps at the ends and four hard edges (Appendix~\ref{app:hardness}). Its ZE partitions are exactly $m$ \emph{positions} $a\in\{0,\dots,m-1\}$, in which the spine members are in states $R^{a}ML^{m-1-a}$ (Lemma~\ref{lem:hchain}). This is a tight instance of bounded $c$ and bounded width in which a ``position'' is selected, so that a cut at any one of arbitrarily many places repairs it (Remark~\ref{rem:hchaincore}).

\paragraph{Thick chains.} Every thin vertex $v$ becomes a \emph{super-vertex} $K_v=\{(v,i):i\in\mathbb Z_p\}$, glued by hard edges, and every thin member $(a,b,c)$ is replaced by the $p^2$ \emph{Latin members} $((a,i),(b,j),(c,i+j))$. These cover every pair between $K_a,K_b$, between $K_b,K_c$ and between $K_a,K_c$ exactly once. All pairs of such a triple of super-vertices are soft; so if $i+j\ne k$, the triple $(a,i),(b,j),(c,k)$ is a soft triangle. The purpose is that a ZE partition cannot split a super-vertex (the glue is hard) and that, for an arbitrary solution of excess at most $\ell$, splitting a super-vertex costs more than $\ell$ hard edits when $p\ge\ell+5$ (Lemma~\ref{lem:hsplit}). Hence the thick chain behaves like the thin chain, while a vertex outside the chain can be attached to a super-vertex at a cost that is independent of the position.

\paragraph{Guests, locks and comparators.} A \emph{guest} is a single new vertex $y$ with a set $W_y$ of super-vertices to which it is soft (through $p/2$ \emph{pad members} $(y,(v,i),(v,k))$ for a perfect matching of $K_v$); all other pairs of $y$ are hard non-edges. A guest can join the cluster of a base vertex $h$ only if the entire cluster of $h$ lies in $W_y$; so the set $W_y$ forbids exactly those positions of the chain at which the cluster of $h$ contains a vertex outside $W_y$ (Lemma~\ref{lem:hprof}). Each chain is divided into $U=1+n(n+1)$ \emph{groups} of two consecutive positions. A \emph{permanent guest} sits in the cluster of a host vertex of each group and, for groups that carry no tile, forbids both positions of the group. For a tile $(a,b)\in S_{ij}$ the chain $A_{ij}$ has a \emph{lock guest} and an \emph{anchor} at the slot group $(a,b)$, and the chain $B_{ij}$ at the slot group $(b,a)$; the two anchors of one tile are joined by a hard edge, a \emph{tunnel}. An uncut tunnel forces the lock guest into the host cluster, hence forbids the group; cutting it (one hard edit, cost one) frees the group. So a chain whose tunnels are all uncut has no position, and a chain with exactly one cut tunnel is in a position of the group of that tunnel; as the same hard edit releases the chains $A_{ij}$ and $B_{ij}$, \emph{it chooses the tile $s_{ij}$ in both orders $(a,b)$ and $(b,a)$ at the price of one edit}. Finally, for consecutive chains of a row (of columns) and a separator group $\alpha<n$, a \emph{comparator guest} $y$ with a member $(z,y,z')$ ($z,z'$ the permanent guests of the two chains) must join exactly one of the two chains, and can do so at the first chain only if its position is before the separator and at the second only if it is after; this is the clause $(\beta\le\alpha)\vee(\beta'>\alpha)$ on the blocks $\beta,\beta'$ of the two chosen tiles, and the $n-1$ clauses together express $\beta\le\beta'$. This is Proposition~\ref{prop:hTR}.

\paragraph{Soundness by counting.} The completeness direction is Proposition~\ref{prop:hTR}. For soundness, an arbitrary solution $S$ of excess at most $\ell=k^2$ is analysed without assuming its form. By Lemma~\ref{lem:hsplit} no super-vertex is split. Partition the pairs into regions $\Pi_c$, one for each of the $2k^2$ chains; every region needs at least one edit beyond one per member (Lemma~\ref{lem:hregion}), because with all tunnels uncut the chain has no position. A pair of two regions is a \emph{cross-anchor pair} (a pair between the anchor sets of the two chains of one cell); only these pairs can serve two regions, and they are hard. If $q$ of them are edited, the excess is at least $\max(q,2k^2-q)\ge k^2$ (Theorem~\ref{thm:hS}); if it equals $k^2$, then $S$ is a set of exactly one tunnel cut per cell plus one edit per member, and Proposition~\ref{prop:hTR} gives a tiling. Thus the instance has $\OPT\le|H|+k^2$ iff the Grid Tiling instance is solvable. In this argument the restricted model of Appendix~\ref{app:tunnels} is not assumed; it is derived.

\paragraph{Where soft triangles are used.} Without soft triangles every instance satisfies the conclusions of Theorem~\ref{thm:sfA}: lock cores have at most two anchors (Corollary~\ref{cor:sfcore}) and tunnel releases cannot realise Theorem~\ref{thm:B} (Theorem~\ref{thm:C}). A single chain with its anchors has a core of unbounded size (Remark~\ref{rem:hchaincore}), so by Lemma~\ref{lem:tmaj} its set of ZE state vectors is closed under no majority operation. The Latin members are the soft triangles: this answers the question of Section~\ref{sec:soft} whether soft triangles suffice for hardness.

\subsection{Computer assistance and what has been tested}\label{sec:hardcomp}
The proof of Theorem~\ref{thm:hard} is \emph{computer-assisted}. Precisely, the following finite statements are established by computation and are used as lemmas; all the other steps (Lemmas~\ref{lem:hblow}, \ref{lem:hsplit}, \ref{lem:hregion}, Proposition~\ref{prop:hTR}, Theorems~\ref{thm:hS} and~\ref{thm:hC}, the pathwidth bound, the pair-maximality) are proved by hand from them.
\begin{enumerate}
\item[(CA1)] \emph{Lemma~\ref{lem:hchain}(a).} The explicit partitions $P_a$ of the thin chain are ZE partitions for all $m\ge7$ and $a$. Each constraint involves at most three consecutive cells; windows at the left end do not depend on $m$, windows at the right end depend only on the distance to the right end, and interior windows are isomorphic, so a check for $7\le m\le60$ and every $a$ (no failure) covers all $m$.
\item[(CA2)] \emph{Lemma~\ref{lem:hchain}(b).} The chain has no other ZE partition. The number $N_m$ of ZE partitions is computed by a transfer-matrix dynamic programme over the cells (six states per interior cell; an explicit integer matrix of size $6\times6$, identical for all interior cells, with characteristic polynomial $x^4(x-1)^2$); this makes $N_m$ an affine function of $m$ for large $m$, and a computation gives $N_m=m$ for $7\le m\le40$. The statement ``the transfer matrix is the same at every interior cell of every $m$'' rests on translation invariance of the member list and was checked at all cells $7,\dots,19$ of $m=26$.
\item[(CA3)] \emph{Lemma~\ref{lem:hprof}.} The partner sets of the host vertices (the clusters of the hosts at the four classes of positions). The statement is read off the formula of the partitions $P_a$; it was in addition checked by computer for all even $m\in[8,60]$ and all groups ($1{,}836$ cases), because the two cap groups have exceptional $e$-labels.
\item[(CA4)] \emph{Counts.} The values $c\le4p+2$ and $u_{\max}\le3$ (Lemma~\ref{lem:hpar}) are obtained by a count that is easily checked by hand from the sets of Lemma~\ref{lem:hprof}; the count was also carried out by computer for $k\le4$.
\end{enumerate}
\emph{Tests that are not part of the proof.} (a) The chain lemma was re-verified by an independent implementation (own validity check and own SAT enumeration of all ZE partitions) for $m=7,\dots,14$ and then for a sample up to $m=100$: always exactly the $m$ partitions $P_a$. (b) The \emph{ZE semantics} of Proposition~\ref{prop:hTR}: for random Grid Tiling instances and every choice of one tile per cell (a sample when there are many), ZE feasibility of the instance with the corresponding tunnels cut agrees with ``the choice is a tiling'' ($k=2$, $n=2,3$; $k=3$, $n=2$; $k=1$, $n=4,5$; thousands of cases in two independent implementations, no mismatch); for $k=n=2$ all $386$ cut sets of size at most $4$ were tested, and exactly the $15$ tilings (one tunnel per cell) are feasible. The comparator was tested on all $100$ pairs of positions of two chains with $m=10$, in the thin model and on the thick instance with $p=8$ ($515$ vertices): feasible exactly when ($a\le3$ or $b\ge6$) and neither $a$ nor $b$ is a permanently forbidden position. A thick instance with $k=2$, $n=2$, $p=8$ ($2{,}916$ vertices, $23{,}176$ members) agreed with the thin model on four choices. (c) \emph{Exact optimum} (all pairs editable; SAT with a cardinality bound on the excess, an UNSAT answer is a proof) on small thick instances only: for $k=1$ (two chains with one tile and its tunnel, $p=8$, $428$ vertices, $|H|=3258$) the optimum exceeds $|H|$ and is at most $|H|+1$ (a certified solution of cost $|H|+1$); the same instance without the tile has $\OPT\ge|H|+2$; and for a comparator instance (two chains, $p=6$, $401$ vertices, $|H|=2349$, $\ell=2$) three YES choices are satisfiable at $|H|+2$ and one NO choice is unsatisfiable at $|H|+2$. \emph{No end-to-end exact solve of an instance with $k\ge2$ was possible} (about $3{,}000$ vertices); for $k=2,3$ the evidence is the tests of the thin model above together with the proof of Theorem~\ref{thm:hS}. (d) Pair-maximality and the validity of the packing were checked on the instances with $(k,n,p)=(1,2,8)$, $(1,2,6)$ to $(1,5,6)$ and $(2,2,10)$ (no induced $P_3$ with three hard pairs). (e) The measured pathwidth upper bounds (an explicit vertex ordering) are $497$ for $k=1$ and $5{,}049$ for $k=2$ at $p=6$ resp.\ $10$ on all components together; they are consistent with Lemma~\ref{lem:hpw}.

% AMENDMENT POINT: findings of any further independent review of Appendix~\ref{app:hardness} are to be recorded in the list above and in the status table of Section~\ref{sec:open}; the appendix is organised so that each lemma can be corrected locally.

\subsection{Consequences for the rest of the paper}
Theorem~\ref{thm:hard} does not contradict any result below: Theorem~\ref{thm:K} needs components of bounded size, but $I(H)$ here has components of size growing with $n$; Theorem~\ref{thm:omega} is fixed-parameter tractable in $(\ell,D,\omega)$ where $D=2c=O(k^2)$ is bounded by a function of $k$, so $\omega$ is unbounded on these instances (the tunnels are hard pairs between far vertices and make the graph $\Gamma'$ of Section~\ref{sec:nl} have unbounded treewidth; compare Proposition~\ref{prop:expander}); Theorem~\ref{thm:sfA} is about soft-triangle-free packings. Two consequences follow from results proved earlier. By Proposition~\ref{prop:rpLOC}, the problem of locating a repacking (Open Problem~\ref{op:LOC}) has no fixed-parameter algorithm unless $\mathrm{FPT}=\mathrm{W[1]}$; and a reduction lemma in the sense of Open Problem~\ref{op:ML} cannot exist unless $\mathrm{FPT}=\mathrm{W[1]}$, since by Theorem~\ref{thm:omega} it would give fixed-parameter tractability in $\ell+w$.
